\documentclass[11pt]{article}
\usepackage{amsmath, amssymb}
\usepackage[margin=1in]{geometry}
\usepackage{booktabs}
\setlength{\parskip}{4pt}
\setlength{\parindent}{0pt}

\newcommand{\py}{\partial_y}
\newcommand{\Shat}{\hat S}
\newcommand{\Hhat}{\hat H}
\newcommand{\col}[1]{\langle #1\rangle}

\title{Moist V1: gross dry stability, mean DSE flux, and gross moist stability,\\
with the equilibrium MSE-flux decomposition}
\date{\today}

\begin{document}
\maketitle

This spells out, from the spec (\texttt{guides/moist\_axisymmetric\_model\_spec}),
the definitions of $\Shat$, the mean DSE flux, and $\Hhat$, evaluates them at the
V1 equilibrium (\texttt{ny}=801, $\Delta t{=}30$, sin$^2$ $y_0{=}0$, slow-drift-gated,
day 1855, $D{=}10^6$), and resolves why the mean flow carries a small
\emph{equatorward} MSE flux while its DSE part is poleward.

\section{Column accounting}

Write the layer thermodynamic and moisture equations of V1 (spec Eqs.~T1, Q1):
\begin{equation}
\partial_t\theta + \frac{d\,\Delta_z}{H}\,\py v = \frac{\theta_E-\theta}{\tau},
\qquad
\partial_t W + \py\!\big[-(2a{-}1)vW - D\,\py W\big] = E_0 - P .
\end{equation}
The column heat capacity that turns a $\theta$-tendency into a column-energy
tendency, and the column dry static energy, are
\begin{equation}
C \equiv \frac{c_p\,\Pi\,\Delta p}{g},
\qquad \col{s}\equiv C\theta .
\end{equation}
With the SS13 Exner factor $\Pi=1/1.6=0.625$, $\Delta p = p_s-p_t = 807$~hPa, and
$c_p=1004~\mathrm{J\,kg^{-1}K^{-1}}$,
\begin{equation}
C = \frac{1004\times 0.625\times 8.07\times10^{4}}{9.81}
  \approx 5.2\times10^{6}~\mathrm{J\,m^{-2}\,K^{-1}} .
\end{equation}
The latent-heating constant is pinned to the \emph{same} $C$ by requiring
$C\Lambda_{\rm conv}=L_v$, i.e.\ $\Lambda_{\rm conv}=L_v/C\approx0.48~\mathrm{K/(kg\,m^{-2}s^{-1})}$
for $L_v=2.5\times10^{6}~\mathrm{J\,kg^{-1}}$.

\section{Gross dry stability and the mean DSE flux}

Multiply the thermodynamic equation by $C$. Since $C$ and the coefficient
$d\Delta_z/H$ are constant, the adiabatic term becomes a flux divergence:
\begin{equation}
\partial_t\col{s} + \py\!\Big[\underbrace{C\,\frac{d\,\Delta_z}{H}}_{\displaystyle \Shat}\,v\Big]
= \frac{C}{\tau}(\theta_E-\theta).
\end{equation}
So the \textbf{gross dry stability} is $\Shat \equiv C\,d\Delta_z/H$ and the
\textbf{mean DSE flux} is $\Shat\,v$. Numerically, with the model's
$d=4000$~m, $\Delta_z=60$~K, $H=1.6\times10^{4}$~m,
\begin{equation}
\frac{d\,\Delta_z}{H}=\frac{4000\times60}{16000}=15~\mathrm{K},
\qquad
\boxed{\;\Shat = C\,\frac{d\Delta_z}{H}=5.2\times10^{6}\times15
      = 7.8\times10^{7}~\mathrm{J\,m^{-2}}\;}
\end{equation}
$\Shat>0$: a poleward $v$ (upper branch, $v>0$ in the NH) exports dry static
energy poleward, cooling the deep tropics, exactly as expected.

\section{Gross moist stability and the MSE budget}

Add $L_v\times$(moisture equation) to $C\times$(thermodynamic equation). Using
$C\Lambda_{\rm conv}=L_v$ the latent heating $\Lambda_{\rm conv}P$ and the
precipitation sink $L_v P$ cancel, leaving a pure column-MSE budget with
$\col{h}=C\theta+L_v W$:
\begin{equation}
\partial_t\col{h}
+\py\!\Big[\underbrace{\Shat v}_{\text{DSE}}
\;\underbrace{-\,L_v(2a{-}1)Wv}_{L_v q}
\;\underbrace{-\,L_v D\,\py W}_{\text{eddy}}\Big]
= \frac{C}{\tau}(\theta_E-\theta) + L_v E_0 .
\end{equation}
The mean (advective) part groups into $v\,\Hhat$ with the \textbf{gross moist
stability}
\begin{equation}
\boxed{\;\Hhat \equiv \Shat - L_v(2a{-}1)\,W\;}
\end{equation}
At the equilibrium equatorial $W\approx50.76~\mathrm{kg\,m^{-2}}$ and $a=0.85$
(so $2a-1=0.7$),
\begin{equation}
\Hhat = 7.8\times10^{7} - (2.5\times10^{6})(0.7)(50.76)
      = 7.8\times10^{7} - 8.88\times10^{7}
      = -1.08\times10^{7}~\mathrm{J\,m^{-2}} .
\end{equation}
$\Hhat<0$: the latent term outweighs the dry stability, the moisture-mode side.
The zero crossing is at $W=\Shat/[L_v(2a{-}1)]=44.6~\mathrm{kg\,m^{-2}}$, which the
column exceeds everywhere, so $\Hhat<0$ across the whole domain.

\section{Why the net mean MSE flux is small and equatorward}

The mean MSE flux $\Shat v - L_v(2a{-}1)Wv = v\Hhat$ is the \emph{difference} of two
large opposing fluxes. At the cell center ($|y|\approx2.9$~Mm, $|v|=0.365~\mathrm{m\,s^{-1}}$):

\begin{center}
\begin{tabular}{lrl}
\toprule
term & MW\,m$^{-1}$ & direction \\
\midrule
DSE $\Shat v$              & $28.5$ & poleward \\
$L_v q$ $-L_v(2a{-}1)Wv$   & $32.4$ & equatorward \\
\midrule
mean MSE $v\Hhat$          & $3.9$  & equatorward (residual) \\
eddy $-L_v D\,\py W$       & $\lesssim 0.15$ & negligible \\
\bottomrule
\end{tabular}
\end{center}

The dry circulation exports $28.5$~MW\,m$^{-1}$ of DSE poleward; the moisture flux,
carried mainly by the equatorward lower branch, returns $32.4$~MW\,m$^{-1}$
equatorward. Their small residual, $3.9$~MW\,m$^{-1}$ up-gradient toward the
ascending ITCZ, is the net MSE flux and the signature of $\Hhat<0$. Plotting only
$v\Hhat$ (as the first diagnostic did) hides that it is the difference of two
$\sim30$~MW\,m$^{-1}$ fluxes; the implied $L_v q$ flux is indeed large.

\section{The overturning is eddy-driven}

The overturning that sets these fluxes follows the eddy-driven leading balance of
the zonal-momentum equation, Coriolis against the eddy momentum flux divergence:
\begin{equation}
\beta y\,v \approx v_d\,\mathcal H(u)\,\mathrm{sgn}(y)\,\py u \equiv \text{EMFD},
\qquad
v \approx \frac{\text{EMFD}}{\beta y}.
\end{equation}
Its magnitude is therefore set by the eddy stress $v_d$ and grows with it. At the
cell center $\beta y\,v = 2.1\times10^{-5}$ balances
$\text{EMFD}=1.8\times10^{-5}~\mathrm{m\,s^{-2}}$ with Rayleigh drag
$\varepsilon_u u = 8.5\times10^{-8}$ two orders of magnitude smaller; the
$\sim13\%$ gap is the nonlinear advection. At the ladder's $v_d=2.5$ this gives
$|v|_{\max}=0.37~\mathrm{m\,s^{-1}}$, so a larger $v_d$ would strengthen the
overturning, its moisture-flux divergence, and the departure of $W$ from local
RCE.

\end{document}
